\section{Folds of finite sequences}

\import{set/finite-sequence.tex}
\import{recursion/natural.tex}

A fold consumes a set-coded finite sequence from left to right.  Its binary
operation is likewise a set-coded function, so the construction composes
directly with Felix's bundled algebraic structures.  Logical natural
recursion supplies the implementation without adding a recursion axiom.

\begin{definition}\label{sequence_fold_step}
    $\sequenceFoldStepOp:
        \Set\to\Set\to\Set\to\Set\to\Set
        =\fun{F,s,n,b:\Set}{
                        \apply{F}{(b,s(n))}}$.
\end{definition}

\begin{definition}\label{sequence_fold_at}
    $\sequenceFoldAt{F}{z}{s}{n}=
        \Apply{\naturalRecOp}{z,
                    \Apply{\sequenceFoldStepOp}{F,s},n}$.
\end{definition}

\begin{definition}\label{sequence_fold}
    $\sequenceFold{F}{z}{s}=
        \sequenceFoldAt{F}{z}{s}{\dom{s}}$.
\end{definition}

\begin{proposition}\label{sequence_fold_at_zero}
    $\sequenceFoldAt{F}{z}{s}{\emptyset}=z$.
\end{proposition}
\begin{proof}
    Follows by
        \obv{sequence_fold_at,natural_recursion_zero}.
\end{proof}

\begin{proposition}\label{sequence_fold_at_successor}
    Suppose $n$ is a natural number.
    Then
        $\sequenceFoldAt{F}{z}{s}{\suc{n}}
            =F((\sequenceFoldAt{F}{z}{s}{n},s(n)))$.
\end{proposition}
\begin{proof}
    Follows by
        \auto{sequence_fold_at,sequence_fold_step,natural_recursion_successor}.
\end{proof}

\begin{proposition}\label{sequence_fold_empty}
    $\sequenceFold{F}{z}{\emptyset}=z$.
\end{proposition}
\begin{proof}
    Follows by
        \auto{sequence_fold,dom_emptyset,sequence_fold_at_zero}.
\end{proof}

\begin{proposition}\label{sequence_fold_at_prefix_extensional}
    Suppose $s$ is a function and $t$ is a function.
    Suppose $n$ is a natural number.
    Suppose $n\subseteq\dom{s}$ and $n\subseteq\dom{t}$.
    Suppose $\restrl{s}{n}=\restrl{t}{n}$.
    Then $\sequenceFoldAt{F}{z}{s}{n}
        =\sequenceFoldAt{F}{z}{t}{n}$.
\end{proposition}
\begin{proof}
    Let $C=\{k\in\naturals\mid
        \text{if $k\subseteq n$, then
            $\sequenceFoldAt{F}{z}{s}{k}
                =\sequenceFoldAt{F}{z}{t}{k}$}\}$.

    $\emptyset\in C$
        by \auto{omega_empty,emptyset,subseteq,sequence_fold_at_zero}.
    Show that for every $k\in C$ we have $\suc{k}\in C$.
    \begin{subproof}
        Fix $k\in C$.
        $\suc{k}$ is a natural number by \obv{omega_successor}.
        Show that if $\suc{k}\subseteq n$, then
            $\sequenceFoldAt{F}{z}{s}{\suc{k}}
                =\sequenceFoldAt{F}{z}{t}{\suc{k}}$.
        \begin{subproof}
            Assume $\suc{k}\subseteq n$.
            $k\subseteq\suc{k}$ by \obv{subseteq_self_suc_intro}.
            Hence $k\subseteq n$ by \obv{subseteq_transitive}.
            Thus $\sequenceFoldAt{F}{z}{s}{k}
                =\sequenceFoldAt{F}{z}{t}{k}$.
            $k\in\suc{k}$ by \obv{suc_intro_self}.
            Hence $k\in n$ by \obv{subseteq}.
            Thus $k\in\dom{s}$ and $k\in\dom{t}$
                by \obv{subseteq}.
            $\apply{\restrl{s}{n}}{k}=s(k)$ and
                $\apply{\restrl{t}{n}}{k}=t(k)$
                by \obv{restrl_of_function_apply}.
            Hence $s(k)=t(k)$.
            Follows by \auto{sequence_fold_at_successor}.
        \end{subproof}
    \end{subproof}
    Hence $C$ is an inductive set by \obv{inductive_set}.
    Hence $\naturals\subseteq C$ by \obv{omega_minimal,subseteq}.
    Thus $n\in C$ by \obv{subseteq}.
    $n\subseteq n$ by \obv{subseteq_refl}.
\end{proof}

\begin{theorem}\label{sequence_fold_append}
    Suppose $s$ is a finite sequence in $A$ and $x\in A$.
    Then
        $\sequenceFold{F}{z}{\sequenceAppend{s}{x}}
            =F((\sequenceFold{F}{z}{s},x))$.
\end{theorem}
\begin{proof}
    Let $n=\dom{s}$.
    Let $t=\sequenceAppend{s}{x}$.
    $n$ is a natural number and $s\in\funs{n}{A}$
        by \auto{finite_sequence_domain}.
    Thus $s$ is a function by \obv{funs_is_function}.
    $t$ is a finite sequence in $A$
        by \auto{sequence_append_finite_sequence}.
    Hence $t$ is a function by
        \auto{finite_sequence_domain,funs_is_function}.
    $\dom{t}=\suc{n}$ by \auto{sequence_append_domain}.
    $n\subseteq\dom{s}$ by \obv{subseteq_refl}.
    $n\subseteq\dom{t}$
        by \obv{subseteq_self_suc_intro}.
    $\restrl{t}{n}=s$ by \auto{sequence_append_restriction}.
    Hence $\restrl{s}{n}=\restrl{t}{n}$
        by \auto{restrl_by_dom}.
    Thus $\sequenceFoldAt{F}{z}{s}{n}
        =\sequenceFoldAt{F}{z}{t}{n}$
        by \obv{sequence_fold_at_prefix_extensional}.
    $t(n)=x$ by \auto{sequence_append_last_value}.
    $\sequenceFoldAt{F}{z}{t}{\suc{n}}
        =F((\sequenceFoldAt{F}{z}{t}{n},t(n)))$
        by \obv{sequence_fold_at_successor}.
    Hence $\sequenceFoldAt{F}{z}{t}{\suc{n}}
        =F((\sequenceFoldAt{F}{z}{s}{n},x))$.
    Follows by \auto{sequence_fold}.
\end{proof}

\begin{theorem}\label{sequence_fold_typed}
    Suppose $s$ is a finite sequence in $A$.
    Suppose $z\in B$ and $F\in\funs{B\times A}{B}$.
    Then $\sequenceFold{F}{z}{s}\in B$.
\end{theorem}
\begin{proof}
    Let $C=\{r\in\finiteSequences{A}\mid
        \sequenceFold{F}{z}{r}\in B\}$.
    $\emptyset\in\finiteSequences{A}$
        by \auto{finite_sequences_iff,emptyset_finite_sequence}.
    Hence $\emptyset\in C$
        by \obv{sequence_fold_empty}.
    Show that for every $r,x$ such that $r\in C$ and $x\in A$ and
        $r$ is a finite sequence in $A$ we have
        $\sequenceAppend{r}{x}\in C$.
    \begin{subproof}
        Fix $r,x$.
        Assume $r\in C$ and $x\in A$ and
            $r$ is a finite sequence in $A$.
        Then $\sequenceFold{F}{z}{r}\in B$.
        Hence $(\sequenceFold{F}{z}{r},x)\in B\times A$
            by \obv{times_tuple_intro}.
        Thus $F((\sequenceFold{F}{z}{r},x))\in B$
            by \obv{funs_type_apply}.
        Hence $\sequenceFold{F}{z}{\sequenceAppend{r}{x}}\in B$
            by \auto{sequence_fold_append}.
        $\sequenceAppend{r}{x}$ is a finite sequence in $A$
            by \obv{sequence_append_finite_sequence}.
        Thus $\sequenceAppend{r}{x}\in\finiteSequences{A}$
            by \obv{finite_sequences_iff}.
    \end{subproof}
    Hence $s\in C$ by \obv{finite_sequence_induction}.
\end{proof}
